Tomando en cuenta los resultados obtenidos durante el desarrollo y validación de la plataforma, se presentan las siguientes recomendaciones y 
consideraciones a tener en cuenta para futuras mejoras o implementaciones:

\begin{itemize}
	\item \textbf{Creación de un módulo de Administración completo.}
		
			Actualmente, la plataforma cuenta con funcionalidades de administración limitadas a la aprobación o rechazo de proyectos
			y cuentas de organizaciones externas por parte de los directores.
			Se recomienda desarrollar un módulo de Administración más robusto que permita a los administradores proyectos, recursos
			y reportes de manera centralizada y eficiente.
	\item \textbf{Gestión de usuarios institucionales por medio del \textit{Single Sign-On} (SSO).}
			
			La implementación actual, lleva a los usuarios institucionales, menciónese a directores y estudiantes, a crear una
			cuenta directamente en la plataforma. 
			Se recomienda integrar el sistema de autenticación de la UVG para que los usuarios puedan iniciar sesión utilizando
			sus credenciales institucionales. Ya que información como el correo electrónico, nombre completo, rol del usuario,
			facultades o carreras relacionadas, entre otros, deben de estar alineados con la información institucional oficial.
	\item \textbf{Pruebas de usabilidad.}
			
			Antes de un posible despliegue de la plataforma, se recomienda realizar más pruebas de usabilidad con una muestra más
			amplia de usuarios finales, para obtener una retroalimentación más representativa y así identificar posibles áreas
			de mejora en la experiencia del usuario.
	\item \textbf{Integración de Horas Beca.}
			
			Se recomienda a la Universidad del Valle de Guatemala (UVG) considerar la integración de un sistema de Horas Beca
			dentro de la plataforma. O bien, que los estudiantes puedan revisar su saldo de horas beca disponibles, así como
			las horas beca que tienen aprobadas y el detalle de las mismas.
\end{itemize}